% Gauss, Werke X/1 (1917), native136–137 upper, printed129–130.
\SourceLeaf{136}\NoSourceHead
\SourceGroup{VI}
\HistoricTitle{[ALGEBRAISCHER LEHRSATZ.]}
\NumberedItem{1}
\HistoricCitation{[Aus Handbuch 21, Bg, S. 76.]}
\noindent\hspace*{\parindent}Eine ganz neue Ansicht der Gleichungen beruhet auf Folgendem. Das
unendliche Planum, welches alle complexen Grössen darstellt, theilt sich scharf
in $n$ Flächenräume, deren jeder Einer der Wurzeln der Gleichung
\[
fx=0
\]
angehört; jeder Scheidungslinie gehört eine Wurzel der Gleichung
\[
f'x=0
\]
an, und in allen Puncten Einer Scheidungslinie ist der imaginäre Theil von
$fx$ constant.

Instructiv wird die Entwicklung der Gleichung
\[
fx=2x^5+5x^4+20x^3+40xx+80x+g
\]
sein, für verschiedene schicklich gewählte Werthe des letzten Gliedes.

Man hat
\[
f'x=10(x^4+2x^3+6xx+8x+8)=10(xx+4)(xx+2x+2)
\]
deren Wurzeln also
\[
\begin{array}{rcl@{\qquad}l}
\alpha&=&+2i&\raisebox{-1.2ex}{$f\alpha=-80-64i+g$}\\
\alpha'&=&-2i&\\[.6ex]
\beta&=&-1+i&\raisebox{-1.2ex}{$f\beta=+42+10i+g$}\\
\beta'&=&-1-i&
\end{array}
\]
\par\vfill\noindent X 1.\hfill17\par
\clearpage
\SourceLeaf{137}\ArithmeticHead{130}{L}{ALGEBRA.\quad NACHLASS.}
\noindent Es ist also
\[
\begin{array}{r@{\quad}l@{\quad}l}
\mathrm{I.}&\text{für }g=+35\tfrac{23}{61}&\operatorname{Mod.}f\alpha=\operatorname{Mod.}f\beta=\tfrac{10}{61}\surd226505\\[.7ex]
\mathrm{II.}&\text{für }g=-49\tfrac79&\operatorname{Ang.}f\alpha=\operatorname{Ang.}f\beta'=180^0+\operatorname{Arc.\ tg.}\tfrac97\\[.7ex]
&&\operatorname{Ang.}f\alpha'=\operatorname{Ang.}f\beta=180^0-\operatorname{Arc.\ tg.}\tfrac97\\[.7ex]
\mathrm{III.}&\text{für }g=-36\tfrac{12}{17}&\operatorname{Ang.}f\alpha+180^0=\operatorname{Ang.}f\beta=\operatorname{Arc.\ tg.}\tfrac{37}{21}\\[.7ex]
&&\operatorname{Ang.}f\alpha'+180^0=\operatorname{Ang.}f\beta'=-\operatorname{Arc.\ tg.}\tfrac{37}{21}.
\end{array}
\]
\noindent\hspace*{\parindent}Einer Entwickelung werth ist die Gleichung
\[
fx=x^5-4x^4+8x^3-8xx+4x+g=0=x(xx-2x+2)^2+g,
\]
wo
\[
\begin{aligned}
f'x&=(5xx-6x+2)(xx-2x+2)\\
&=5\left(x-\frac{3+i}{5}\right)\left(x-\frac{3-i}{5}\right)(x-(1+i))(x-(1-i)).
\end{aligned}
\]
Namentlich in Beziehung auf die Linie, wo
\[
f(a+bi)=f(a-bi).
\]
\ClosingRule
